%%% BEGIN fig:method-nesting
\begin{figure}[htbp]
\centering
\begin{tikzpicture}[fd,
  fb lab/.style={font=\footnotesize},
  fb ellA/.style={rotate around={22:(-0.95,0)}},
  fb ellB/.style={rotate around={-22:(0.95,0)}},
  fb step/.style={draw=gray!40!black,line width=0.75pt,-{Stealth[length=5pt,width=4.2pt]}},
  fb stepl/.style={font=\footnotesize,inner sep=1.5pt,fill=white,rounded corners=1.5pt}]
  % ---------------- Venn-like diagram of the pinch sets (schematic) ----------------
  \draw[fill=gray!10,draw=gray!55!black,line width=0.7pt,rounded corners=10pt] (-3.6,-2.35) rectangle (3.6,2.35);
  \fill[blue!12,fb ellA] (-0.95,0) ellipse[x radius=2.05,y radius=1.08];
  \fill[orange!17,fb ellB] (0.95,0) ellipse[x radius=2.05,y radius=1.08];
  % pi_S = pi_A cap pi_B: the lens common to both ellipses
  \begin{scope}
    \clip[fb ellA] (-0.95,0) ellipse[x radius=2.05,y radius=1.08];
    \fill[green!22,fb ellB] (0.95,0) ellipse[x radius=2.05,y radius=1.08];
  \end{scope}
  \draw[blue!55!black,line width=0.7pt,fb ellA] (-0.95,0) ellipse[x radius=2.05,y radius=1.08];
  \draw[orange!75!black,line width=0.7pt,fb ellB] (0.95,0) ellipse[x radius=2.05,y radius=1.08];
  \node[font=\small,text=green!30!black] at (0,0.22) {$\pi_S$};
  \node[font=\small,text=blue!50!black] at (-2.05,-0.25) {$\pi_A$};
  \node[font=\small,text=orange!60!black] at (2.05,-0.25) {$\pi_B$};
  \node[font=\small,anchor=north west] at (-3.5,2.25) {$\pi_H=\Omega$};
  \node[fb lab,text=gray!30!black,anchor=south] at (0,-2.3) {$\pi_S=\pi_A\cap\pi_B$（\cref{lem:method-meets}）};
  % order of the approximations: smallest region first
  \draw[fb step] (0.1,1.95) node[fb stepl,anchor=south] {\textcircled{\scriptsize 1}\,$T_S$} to[out=-90,in=90] (0.04,0.62);
  \draw[fb step] (-2.9,-1.65) node[fb stepl,anchor=north] {\textcircled{\scriptsize 2}\,$T_A$} to[out=60,in=-150] (-1.85,-0.75);
  \draw[fb step] (2.9,-1.65) node[fb stepl,anchor=north] {\textcircled{\scriptsize 2}\,$T_B$} to[out=120,in=-30] (1.85,-0.75);
  \draw[fb step] (2.75,1.85) node[fb stepl,anchor=south west,xshift=-0.15cm] {\textcircled{\scriptsize 3}\,$T_H$} to[out=-120,in=60] (2.4,1.2);
  % ---------------- subtraction terms ----------------
  \begin{scope}[xshift=4.15cm]
    \node[anchor=north west,align=left,font=\small] (F) at (0,2.2) {%
      $C_S\Gamma=T_S\,\Gamma$\\[4pt]
      $C_A\Gamma=T_A(1-T_S)\,\Gamma$，\quad$C_B\Gamma=T_B(1-T_S)\,\Gamma$\\[4pt]
      $C_H\Gamma=T_H\big(\Gamma-C_S\Gamma-C_A\Gamma-C_B\Gamma\big)$\\[4pt]
      $\Gamma=\sum_{R}C_R\Gamma+r$};
    \node[anchor=north west,align=left,font=\footnotesize,text width=6.6cm,draw=red!50!black,
          line width=0.5pt,rounded corners=2pt,fill=red!3,inner sep=4pt] at ($(F.south west)+(0.05,-0.25)$)
      {只有分组表达式被压低（\cref{ex:method-naive-induction}）：单个链项一般不可积（例如离壳三角图中的 $T_{C_1}\Gamma$ 对任何 $\varepsilon$ 都不绝对可积），被压低并满足 \mbox{(CS1)--(CS3)} 的是分组表达式 $C_R\Gamma$ 与余项 $r$。};
  \end{scope}
\end{tikzpicture}
\caption{区域的偏序与减除（示意；钳夹集画成区域）。单圈顶点的钳夹集彼此嵌套：软钳夹集 $\pi_S$（绿色）同时含于两个共线钳夹集 $\pi_A$（蓝）与 $\pi_B$（橙），并且恰为二者之交 $\pi_S=\pi_A\cap\pi_B$，即逐线交 $A\wedge B=S$；所有钳夹集都含于硬区 $\pi_H=\Omega$。
减除自最小区域（软）开始，依次为 \textcircled{\scriptsize 1}~$T_S$、\textcircled{\scriptsize 2}~$T_A,T_B$、\textcircled{\scriptsize 3}~$T_H$：较大区域的近似子作用于减去较小区域贡献之后的被积函数，$C_R\Gamma\equiv T_R\big(\Gamma-\sum_{R'<R}C_{R'}\Gamma\big)$（\cref{def:method-subtraction}）。
单个链项一般不可积，只有分组表达式满足 \mbox{(CS1)--(CS3)}。}
\label{fig:method-nesting}
\end{figure}
%%% END fig:method-nesting
